\section{Conclusion}

In this work, we showed that iterative deployment of LLMs, where each generation
is fine-tuned on curated traces from previous deployments, improves planning
capabilities. Our experimental results on classical planning domains show that
models trained under this paradigm more than doubled their performance within
five generations, with evidence of out-of-distribution generalization to
longer-horizon plans.

Iterative deployment can be an effective alternative to RL to improve reasoning
capabilities of LLMs. We proved theoretically that our method is equivalent to a
special case of REINFORCE with binary rewards. Unlike RL fine-tuning, our
approach does not rely on carefully designed reward functions. Instead, it uses
external postprocessing tools to validate/curate LLM answers as an implicit
signal.  However, the absence of an internal reward function brings concern
about potential biases and unintended effects caused during training. The use of
external validators, that might be outside our control, can be damaging to AI
safety.

Overall, our results suggest that iterative deployment is a viable method for
self-improving LLMs. In future work, we plan to study how model collapsing
\citep{shumailov-et-al-nature2024} can affect iterative deployment, and also to
develop theoretical results to connect our method with RL fine-tuning more
explicitly.
